% nexus-icons.tex — Jeu d'icônes propriétaires TikZ pour le manuel Nexus Réussite
% Trait 0.8pt uniforme, boîte 1em, monochromes (couleur du contexte).
% Remplace tous les \ding de pifont.

\RequirePackage{tikz}
\usetikzlibrary{calc}

% --- Dimensions communes ---
\newlength{\nxIconSize}
\setlength{\nxIconSize}{1em}

% ============================================================
%  LOSANGE DE PARCOURS (remplace \ding{117})
% ============================================================
\newcommand{\nxLosange}[1][chapcolor]{%
  \tikz[baseline=-0.15em]{%
    \fill[#1] (0,0) -- (0.3em,0.45em) -- (0.6em,0) -- (0.3em,-0.45em) -- cycle;
    \draw[encre,line width=0.4pt,opacity=0.3] (0,0) -- (0.3em,0.45em) -- (0.6em,0) -- (0.3em,-0.45em) -- cycle;
  }%
}

% Parcours redessinés
\renewcommand{\parcoursUn}{\nxLosange}
\renewcommand{\parcoursDeux}{\nxLosange\kern0.15em\nxLosange}
\renewcommand{\parcoursTrois}{\nxLosange\kern0.15em\nxLosange\kern0.15em\nxLosange}

% ============================================================
%  \icnObjectif — cible (« je sais… »)
% ============================================================
\newcommand{\icnObjectif}[1][nxBleu]{%
  \tikz[baseline=-0.15em]{%
    \draw[#1,line width=0.8pt] (0,0) circle (0.45em);
    \draw[#1,line width=0.8pt] (0,0) circle (0.25em);
    \fill[#1] (0,0) circle (0.08em);
  }%
}

% ============================================================
%  \icnMethode — boussole
% ============================================================
\newcommand{\icnMethode}[1][nxOr]{%
  \tikz[baseline=-0.15em]{%
    \draw[#1,line width=0.8pt] (0,0) circle (0.45em);
    \draw[#1,line width=0.8pt] (0,-0.35em) -- (0.12em,0.05em) -- (0,0.35em) -- (-0.12em,0.05em) -- cycle;
    \fill[#1] (0,0.05em) -- (0.12em,0.05em) -- (0,0.35em) -- cycle;
  }%
}

% ============================================================
%  \icnErreur — triangle arrondi + !
% ============================================================
\newcommand{\icnErreur}[1][nxRouge]{%
  \tikz[baseline=-0.15em]{%
    \draw[#1,line width=0.8pt,rounded corners=0.8pt]
      (-0.4em,-0.35em) -- (0,0.45em) -- (0.4em,-0.35em) -- cycle;
    \draw[#1,line width=0.8pt] (0,-0.05em) -- (0,0.2em);
    \fill[#1] (0,-0.2em) circle (0.05em);
  }%
}

% ============================================================
%  \icnCle — clé (coups de pouce)
% ============================================================
\newcommand{\icnCle}[1][nxGris]{%
  \tikz[baseline=-0.15em]{%
    \draw[#1,line width=0.8pt] (0.15em,0) circle (0.22em);
    \draw[#1,line width=0.8pt] (-0.07em,0) -- (-0.5em,0);
    \draw[#1,line width=0.8pt] (-0.35em,0) -- (-0.35em,-0.15em);
    \draw[#1,line width=0.8pt] (-0.5em,0) -- (-0.5em,-0.15em);
  }%
}

% ============================================================
%  \icnCheck — coche (corrigés)
% ============================================================
\newcommand{\icnCheck}[1][nxVert]{%
  \tikz[baseline=-0.15em]{%
    \draw[#1,line width=0.8pt,line cap=round]
      (-0.3em,0) -- (-0.1em,-0.3em) -- (0.35em,0.3em);
  }%
}

% ============================================================
%  \icnOral — bulle de parole
% ============================================================
\newcommand{\icnOral}[1][nxBleu]{%
  \tikz[baseline=-0.15em]{%
    \draw[#1,line width=0.8pt,rounded corners=1.5pt]
      (-0.4em,-0.15em) rectangle (0.4em,0.35em);
    \draw[#1,line width=0.8pt]
      (-0.1em,-0.15em) -- (-0.15em,-0.35em) -- (0.1em,-0.15em);
  }%
}

% ============================================================
%  \icnPython — chevrons ›››
% ============================================================
\newcommand{\icnPython}[1][nxBleu]{%
  \tikz[baseline=-0.15em]{%
    \foreach \x in {0,0.22em,0.44em} {
      \draw[#1,line width=0.8pt] (\x,-0.25em) -- (\x+0.15em,0) -- (\x,0.25em);
    }
  }%
}

% ============================================================
%  \icnChrono — cadran (durées)
% ============================================================
\newcommand{\icnChrono}[1][nxGris]{%
  \tikz[baseline=-0.15em]{%
    \draw[#1,line width=0.8pt] (0,0) circle (0.4em);
    \draw[#1,line width=0.8pt] (0,0) -- (0,0.25em);
    \draw[#1,line width=0.8pt] (0,0) -- (0.2em,0.1em);
    \draw[#1,line width=0.8pt] (-0.05em,0.45em) -- (0.05em,0.45em);
  }%
}

% ============================================================
%  \icnEtoile — étoile 4 branches fines (★)
% ============================================================
\newcommand{\icnEtoile}[1][nxBleu]{%
  \tikz[baseline=-0.15em]{%
    \foreach \a in {0,90,180,270} {
      \draw[#1,line width=0.8pt] (0,0) -- (\a:0.4em);
    }
    \foreach \a in {45,135,225,315} {
      \draw[#1,line width=0.8pt] (0,0) -- (\a:0.2em);
    }
  }%
}

% ============================================================
%  \icnCarte — nœuds reliés (carte du chapitre)
% ============================================================
\newcommand{\icnCarte}[1][nxGris]{%
  \tikz[baseline=-0.15em]{%
    \fill[#1] (-0.3em,0.2em) circle (0.07em);
    \fill[#1] (0,0.35em) circle (0.07em);
    \fill[#1] (0.3em,0.2em) circle (0.07em);
    \fill[#1] (-0.15em,-0.2em) circle (0.07em);
    \fill[#1] (0.15em,-0.2em) circle (0.07em);
    \draw[#1,line width=0.8pt] (-0.3em,0.2em) -- (0,0.35em) -- (0.3em,0.2em);
    \draw[#1,line width=0.8pt] (0,0.35em) -- (-0.15em,-0.2em);
    \draw[#1,line width=0.8pt] (0,0.35em) -- (0.15em,-0.2em);
    \draw[#1,line width=0.8pt] (-0.15em,-0.2em) -- (0.15em,-0.2em);
  }%
}

% ============================================================
%  Mise à jour des commandes Oral et Python
% ============================================================
\renewcommand{\pictoOral}{\icnOral[chapcolor]~\textcolor{chapcolor}{Oral}}
\renewcommand{\pictoPython}{\icnPython[chapcolor]~\textcolor{chapcolor}{\texttt{Python}}}
